\section{Algorithm Flowcharts and Solver Workflows}
\label{sec:algorithms}

\subsection{Canonical Preprocessing Workflow}

\begin{figure}[H]
\centering
\begin{tikzpicture}[node distance=7mm and 10mm]
  \node[flowbox] (start) {Input \texttt{HybridPowerSystem}};
  \node[flowbox, below=of start] (proj) {Project rich models to canonical equivalents};
  \node[flowbox, below=of proj] (agg) {Aggregate per-bus generation and demand};
  \node[flowbox, below=of agg] (mat) {Build \(Y_{bus}\), \(G_{dc}\), ZIP tables};
  \node[flowbox, below=of mat] (dispatch) {Dispatch to PF / OPF / SC solver path};
  \draw[flowline] (start) -- (proj);
  \draw[flowline] (proj) -- (agg);
  \draw[flowline] (agg) -- (mat);
  \draw[flowline] (mat) -- (dispatch);
\end{tikzpicture}
\caption{Shared preprocessing path used by the main PF, native OPF, parity OPF, and SC modules.}
\end{figure}

\subsection{Hybrid Newton Power Flow}

\begin{figure}[H]
\centering
\begin{tikzpicture}[node distance=7mm and 12mm]
  \node[flowbox] (init) {Initialize \(v,\theta,v^{dc}\) from case or warm start};
  \node[flowbox, below=of init] (ctx) {Build bus-type maps and Jacobian context};
  \node[flowbox, below=of ctx] (eval) {Evaluate residuals and Jacobian};
  \node[decision, below=of eval] (conv) {Residual below tolerance?};
  \node[flowbox, right=30mm of conv] (done) {Populate results and derived flows};
  \node[flowbox, below=of conv] (linsolve) {Factorize / regularize and solve \(J\Delta x = F\)};
  \node[flowbox, below=of linsolve] (clip) {Clip step and run backtracking line search};
  \node[flowbox, below=of clip] (switch) {Apply PV/PQ and converter-mode switching if enabled};
  \draw[flowline] (init) -- (ctx);
  \draw[flowline] (ctx) -- (eval);
  \draw[flowline] (eval) -- (conv);
  \draw[flowline] (conv) -- node[above]{yes} (done);
  \draw[flowline] (conv) -- node[right]{no} (linsolve);
  \draw[flowline] (linsolve) -- (clip);
  \draw[flowline] (clip) -- (switch);
  \draw[flowline] (switch.west) |- ++(-18mm,0) |- (eval.west);
\end{tikzpicture}
\caption{Main Newton-Raphson workflow for hybrid AC/DC power flow.}
\end{figure}

\subsection{AC OPF Path Selection}

\begin{figure}[H]
\centering
\resizebox{\linewidth}{!}{%
\begin{tikzpicture}[
  node distance=12mm and 22mm,
  every node/.style={font=\small}
]
  \node[flowwide] (start) {Start \texttt{solve\_ac\_opf}};
  \node[decision, below=of start, text width=3.1cm] (pd) {\texttt{enable\_primal\_dual}?};
  \node[flowwide, left=60mm of pd, text width=4.8cm] (fast) {Fast economic-dispatch + AC PF fallback};
  \node[decision, below=of pd, text width=3.1cm] (parity) {\texttt{use\_parity\_ipm}?};
  \node[flowwide, left=60mm of parity, text width=4.8cm] (native) {Native reduced-space primal-dual OPF};
  \node[flowwide, below=of parity, text width=5.0cm] (parsolve) {Build parity NLP and solve parity IPM};
  \node[decision, below=of parsolve, text width=2.8cm] (ok) {Converged?};
  \node[flowwide, below=of ok, text width=5.2cm] (fallback) {Fallback to native or fast AC-only path (if allowed)};
  \node[flowwide, right=54mm of ok, text width=4.3cm] (done) {Return \texttt{ACOPFResult}};
  \draw[flowline] (start) -- (pd);
  \draw[flowline] (pd) -- node[above]{no} (fast);
  \draw[flowline] (pd) -- node[right]{yes} (parity);
  \draw[flowline] (parity) -- node[above]{no} (native);
  \draw[flowline] (parity) -- node[right]{yes} (parsolve);
  \draw[flowline] (parsolve) -- (ok);
  \draw[flowline] (ok) -- node[above]{yes} (done);
  \draw[flowline] (ok) -- node[right]{no} (fallback);
  \draw[flowline] (fast.east) -- ++(12mm,0) |- (done.north west);
  \draw[flowline] (native.east) -- ++(12mm,0) |- (done.west);
  \draw[flowline] (fallback.east) -- ++(12mm,0) |- (done.south west);
\end{tikzpicture}
}
\caption{Current top-level solver-selection logic for AC OPF.}
\end{figure}

\subsection{Hybrid Carbon-Analysis Workflow}

\begin{figure}[H]
\centering
\begin{tikzpicture}[node distance=7mm and 12mm]
  \node[flowbox] (proj) {Project rich models to canonical AC/DC network};
  \node[flowbox, below=of proj] (pf) {Read converged PF state: AC flows, DC voltages, VSC / DCDC transfers, router ports};
  \node[flowbox, below=of pf] (graph) {Build directed hybrid carrier graph with AC/DC branches, converters, and router-internal carriers};
  \node[flowbox, below=of graph] (norm) {Normalize explicit carbon sources, storage sinks, and unified AC/DC loads};
  \node[flowbox, below=of norm] (trace) {Run upstream proportional tracing for loads and carrier losses};
  \node[flowbox, below=of trace] (matrix) {Assemble hybrid matrix system and solve bus carbon intensities};
  \node[flowbox, below=of matrix] (report) {Export AC/DC load, branch, bus, storage, VSC, DCDC, and EnergyRouter carbon reports};
  \draw[flowline] (proj) -- (pf);
  \draw[flowline] (pf) -- (graph);
  \draw[flowline] (graph) -- (norm);
  \draw[flowline] (norm) -- (trace);
  \draw[flowline] (trace) -- (matrix);
  \draw[flowline] (matrix) -- (report);
\end{tikzpicture}
\caption{Current hybrid carbon-analysis workflow.}
\end{figure}

\subsection{Distribution PF Workflow}

\begin{figure}[H]
\centering
\begin{tikzpicture}[node distance=7mm and 12mm]
  \node[flowbox] (tree) {Check radial feeder and build BFS tree};
  \node[flowbox, below=of tree] (inj) {Assemble net demand \(S^{net}\) from direct AC tables};
  \node[flowbox, below=of inj] (back) {Backward sweep: currents from leaves to root};
  \node[flowbox, below=of back] (forw) {Forward sweep: voltages from root to leaves};
  \node[decision, below=of forw] (res) {Voltage-change residual below tolerance?};
  \node[flowbox, right=32mm of res] (report) {Compute branch powers and losses};
  \draw[flowline] (tree) -- (inj);
  \draw[flowline] (inj) -- (back);
  \draw[flowline] (back) -- (forw);
  \draw[flowline] (forw) -- (res);
  \draw[flowline] (res) -- node[above]{yes} (report);
  \draw[flowline] (res.west) -- ++(-16mm,0) |- node[left]{no} (back.west);
\end{tikzpicture}
\caption{Backward-forward sweep workflow for radial distribution PF.}
\end{figure}

\subsection{Short-Circuit Workflow}

\begin{figure}[H]
\centering
\begin{tikzpicture}[node distance=7mm and 12mm]
  \node[flowbox] (proj) {Project to canonical AC model};
  \node[flowbox, below=of proj] (ybus) {Assemble fault admittance with branch, generator, motor, and external-grid terms};
  \node[flowbox, below=of ybus] (fact) {Factorize \(Y_{fault}\)};
  \node[flowbox, below=of fact] (solve) {For each fault bus, solve \(Y_{fault}z=e_k\)};
  \node[flowbox, below=of solve] (calc) {Recover \(Z_{kk}\), fault current, and short-circuit power};
  \draw[flowline] (proj) -- (ybus);
  \draw[flowline] (ybus) -- (fact);
  \draw[flowline] (fact) -- (solve);
  \draw[flowline] (solve) -- (calc);
\end{tikzpicture}
\caption{Primary short-circuit computation path in the current module.}
\end{figure}

\subsection{ONR Workflow}

\begin{figure}[H]
\centering
\begin{tikzpicture}[node distance=7mm and 12mm]
  \node[flowbox] (net) {Build direct AC net-injection vectors};
  \node[flowbox, below=of net] (milp) {Assemble LinDistFlow MILP};
  \node[flowbox, below=of milp] (bc) {Solve with native branch-and-cut};
  \node[flowbox, below=of bc] (snap) {Snap branch-status variables to a spanning tree};
  \node[flowbox, below=of snap] (verify) {Run Newton PF on selected topology};
  \draw[flowline] (net) -- (milp);
  \draw[flowline] (milp) -- (bc);
  \draw[flowline] (bc) -- (snap);
  \draw[flowline] (snap) -- (verify);
\end{tikzpicture}
\caption{Optimal network reconfiguration workflow.}
\end{figure}

\subsection{Solver-Family Summary}

\begin{longtable}{L{0.22\linewidth}L{0.29\linewidth}L{0.39\linewidth}}
\toprule
Workflow & Numerical core & Main reason to use it \\
\midrule
\endhead
Hybrid Newton PF & sparse Newton + line search + regularization & highest-fidelity steady-state solve in the current module. \\
DC-only PF & dense Newton on DC voltages & fast solution for purely resistive DC subnetworks. \\
FDPF & sparse decoupled angle/voltage solves & faster AC approximation when converter coupling is irrelevant. \\
Adaptive PF & island detection + per-island Newton solves & robust handling of disconnected or partially energized systems. \\
Distributed slack & post-PF mismatch redistribution & studies with multi-source slack participation. \\
Native AC OPF & custom KKT system & detailed nonlinear dispatch with direct generator/converter variables. \\
Parity OPF & parity-form NLP + dedicated IPM & full-space formulation with generator, VSC, DCDC, load-shedding, and DER decision variables; analytical Jacobian and Hessian including \(I^2R\) DCDC loss terms. \\
DC OPF & LP/QP on \(B\)-\(\theta\) network & fast linearized economic dispatch with LMP extraction; supports piecewise-linear and quadratic cost models. \\
Reactive power opt.\ (RPO) & MINLP with 4-phase OA & discrete tap/shunt optimization with AC OPF as inner NLP. \\
LCQP solver & Mehrotra predictor-corrector IPM & native convex QP/LP solver used as DC OPF backbone and RPO sub-solver. \\
Dual simplex & two-phase simplex with basis caching & LP for DC OPF and branch-and-cut MILP sub-problems; warm-start support. \\
Hybrid carbon analysis & graph tracing + dense QR on hybrid carrier graph & post-PF carbon propagation across AC branches, DC branches, storage charge/discharge states, VSC links, DCDC links, and EnergyRouter carriers. \\
BFS DPF & tree-based current/voltage sweep & radial feeder studies where a full Jacobian solve is unnecessary. \\
ONR branch-and-cut & MILP B\&C with MIR cuts & topology optimization over switch states. \\
Z-bus SC & sparse factorization of fault admittance & per-bus fault-level studies on the canonical AC model. \\
Annual production sim & hierarchical L0/L2/L3 decomposition & year-ahead scheduling with UC + OPF/PF replay. \\
Lifecycle simulation & multi-year degradation + stratified carbon & techno-economic planning with battery replacement and error bounds. \\
\bottomrule
\end{longtable}

\subsection{Code-Aligned Pseudocode (All Primary Algorithms)}
\label{sec:algo-pseudocode}

\paragraph{Algorithm A0: Runtime API Dispatch}
\begin{algorithmic}[1]
\Procedure{\texttt{run\_module}}{$sys, task, opt$}
  \If{$task=\texttt{"power\_flow"}$}
    \State \Return \texttt{solve\_power\_flow}$(sys, opt.pf)$
  \ElsIf{$task=\texttt{"dc\_power\_flow"}$}
    \State \Return \texttt{solve\_dc\_power\_flow}$(sys, opt.pf)$
  \ElsIf{$task=\texttt{"power\_flow\_adaptive"}$}
    \State \Return \texttt{solve\_power\_flow\_adaptive}$(sys, opt.pf)$
  \ElsIf{$task=\texttt{"opf"}$}
    \State \Return \texttt{solve\_ac\_opf}$(sys, opt.opf)$
  \ElsIf{$task=\texttt{"short\_circuit"}$}
    \State \Return \texttt{compute\_short\_circuit}$(sys, opt.sc)$
  \ElsIf{$task=\texttt{"network\_reconfiguration"}$}
    \State \Return \texttt{solve\_optimal\_reconfiguration}$(sys, opt.onr)$
  \ElsIf{$task=\texttt{"distribution\_pf"}$}
    \State \Return \texttt{solve\_distribution\_pf}$(sys, opt.dpf)$
  \ElsIf{$task=\texttt{"carbon\_analysis"}$}
    \State $pf \gets \texttt{solve\_power\_flow}(sys, opt.pf)$
    \State \Return \texttt{compute\_carbon\_analysis}$(sys, pf, opt.carbon)$
  \Else
    \State \Return \texttt{UnsupportedTask}
  \EndIf
\EndProcedure
\end{algorithmic}

\paragraph{Algorithm A1: Canonical Projection + SolverData Build}
\begin{algorithmic}[1]
\Procedure{\texttt{make\_solver\_data}}{$sys, loss\_model$}
  \State $projected \gets \texttt{project\_to\_canonical\_models}(sys)$
  \Statex \Comment{Includes rich-model expansion and optional zero-\(Z\) bus merging}
  \State $data \gets$ copy canonical tables from $projected$
  \State $data.loss\_model \gets loss\_model$
  \State Apply external-grid voltage/slack settings
  \State Aggregate generation and demand
  \State Build \(Y_{bus}\), \(G_{dc}\), and ZIP tables
  \State \Return $data$
\EndProcedure
\end{algorithmic}

\paragraph{Algorithm A2: \texttt{solve\_power\_flow} (Top Level)}
\begin{algorithmic}[1]
\Procedure{\texttt{solve\_power\_flow}}{$sys, opt$}
  \State $data \gets$ cached \texttt{SolverData} for $(sys, opt.loss\_model)$
  \State Apply ZIP option weights to $data$
  \State $result \gets \texttt{NewtonSolver.solve}(data, opt)$
  \State Populate branch/converter/DCDC/router derived results
  \If{$data.bus\_merge\_map$ exists}
    \State Unproject \(v_m,\theta\) back to the original AC-bus count
  \EndIf
  \State \Return $result$
\EndProcedure
\end{algorithmic}

\paragraph{Algorithm A3: Hybrid Newton PF Core}
\begin{algorithmic}[1]
\Procedure{\texttt{NewtonSolver.solve}}{$data, opt, init$}
  \State Initialize \(v,\theta,v^{\qual{dc}}\), slack sets, and control modes
  \For{$outer=1$ to \texttt{max\_outer}}
    \State Build Jacobian context and sparse pattern caches
    \For{$iter=1$ to \texttt{opt.max\_iter}}
      \State Evaluate residual and Jacobian
      \If{$\|F\|_{\infty} < \texttt{tol}$} \State \textbf{break} \EndIf
      \State Solve \(J\Delta x = F\) with regularization fallback
      \State Clip step, then run backtracking line search
      \If{no acceptable step} \State mark failure; \textbf{break} \EndIf
      \State Update state with accepted step
      \State Apply PV\(\rightarrow\)PQ and converter-mode switching rules
    \EndFor
    \If{converged and PQ\(\rightarrow\)PV restoration is enabled}
      \State Restore eligible PV buses and continue outer loop
    \Else
      \State \textbf{break}
    \EndIf
  \EndFor
  \State Export \(v,\theta,v^{\qual{dc}}\), residual, and status
  \State \Return result
\EndProcedure
\end{algorithmic}

\paragraph{Algorithm A4: DC-Only Newton PF}
\begin{algorithmic}[1]
\Procedure{\texttt{DCSolver.solve}}{$data, opt, init$}
  \State Initialize \(v^{\qual{dc}}\), detect DC slack, build reduced index set
  \For{$iter=1$ to \texttt{max\_iter}}
    \State Assemble \(P^{\qual{dc,spec}}\) from bus/load/storage/gen/VSC/DCDC/router
    \State Compute \(P^{\qual{dc,calc}} = v^{\qual{dc}} \odot (G_{dc}v^{\qual{dc}})\)
    \State Form mismatch on non-slack rows
    \If{$\|mismatch\|_{\infty} < \texttt{tol}$} \State converged; \textbf{break} \EndIf
    \State Build dense reduced Jacobian and solve for \(\Delta v^{\qual{dc}}\)
    \State Apply positivity-safe line search
  \EndFor
  \State \Return \(v^{\qual{dc}}\), residual, convergence flag
\EndProcedure
\end{algorithmic}

\paragraph{Algorithm A5: Fast-Decoupled PF (FDPF)}
\begin{algorithmic}[1]
\Procedure{\texttt{FDPFSolver.solve}}{$data, opt$}
  \State Build reduced \(B'\) and \(B''\); factorize once
  \State Initialize \(v,\theta\)
  \For{$iter=1$ to \texttt{max\_iter}}
    \State Compute \((P^{\qual{calc}},Q^{\qual{calc}})\) and ZIP-adjusted specs
    \State Form \(P\)- and \(Q\)-mismatch vectors
    \If{both mismatch norms \(<\texttt{tol}\)} \State converged; \textbf{break} \EndIf
    \State Solve \(B' \Delta\theta = -\Delta P\), \(B''\Delta v = -\Delta Q\)
    \State Update non-slack \(\theta\) and PQ-bus \(v\)
  \EndFor
  \State \Return \(v,\theta\), residual, convergence flag
\EndProcedure
\end{algorithmic}

\paragraph{Algorithm A6: Adaptive/Islanded PF}
\begin{algorithmic}[1]
\Procedure{\texttt{AdaptiveSolver.solve}}{$sys, opt$}
  \State $islands \gets \texttt{detect\_islands}(sys)$
  \If{$islands$ is empty}
    \State Solve full system once with Newton PF and return
  \EndIf
  \State Initialize global output vectors
  \For{each island}
    \If{island has no generation}
      \State Set island voltages to zero and continue
    \EndIf
    \State Auto-select swing bus if required
    \State Extract subsystem and run Newton PF (with one flat-start retry)
    \State Unproject merged AC results if bus merging occurred
    \State Map local results back to global vectors
  \EndFor
  \State Aggregate iteration count and residual
  \State \Return global adaptive result
\EndProcedure
\end{algorithmic}

\paragraph{Algorithm A7: Distributed Slack}
\begin{algorithmic}[1]
\Procedure{\texttt{solve\_distributed\_slack}}{$sys, slack\_cfg, opt$}
  \State Build/validate participation configuration
  \State Solve base PF
  \If{base PF not converged} \State \Return base result \EndIf
  \State Compute total active-power mismatch at participating buses
  \State Allocate mismatch by participation factors and optional limits
  \State Apply allocation and export adjusted injections/summary
  \State \Return distributed-slack result
\EndProcedure
\end{algorithmic}

\paragraph{Algorithm A8: AC OPF Top-Level Router}
\begin{algorithmic}[1]
\Procedure{\texttt{solve\_ac\_opf}}{$sys, opt$}
  \State Normalize default options
  \State Detect hybrid case (\(DC\) buses or VSC converters present)
  \If{hybrid case and \texttt{enable\_primal\_dual=false}}
    \State Force \texttt{enable\_primal\_dual=true}, \texttt{use\_parity\_ipm=true}
  \EndIf
  \If{\texttt{enable\_primal\_dual=true} and \texttt{use\_parity\_ipm=true}}
    \State Solve parity IPM path
    \If{failed and \texttt{allow\_fallback}} \State retry native path \EndIf
    \State \Return result
  \EndIf
  \If{\texttt{enable\_primal\_dual=false}}
    \State \Return fast dispatch + AC PF fallback path
  \EndIf
  \State \Return native primal-dual OPF path
\EndProcedure
\end{algorithmic}

\paragraph{Algorithm A9: Fast Dispatch + AC PF Fallback (OPF)}
\begin{algorithmic}[1]
\Procedure{\texttt{try\_dispatch\_pf\_fallback}}{$ac\_sys, opt$}
  \State Try direct PF projection using current generator settings
  \If{PF converged}
    \State Back-calculate bus injections and split by generator limits
    \State Compute objective and return converged OPF result
  \EndIf
  \State Solve clipped economic dispatch (\(\lambda\)-search)
  \State Rerun AC PF projection with dispatched \(P_g\)
  \If{PF converged} \State \Return converged result \EndIf
  \State \Return not-converged result
\EndProcedure
\end{algorithmic}

\paragraph{Algorithm A10: Native Primal-Dual AC/DC OPF}
\begin{algorithmic}[1]
\Procedure{\texttt{solve\_native\_primal\_dual}}{$sys, opt$}
  \State Build \(x\)-map, bounds, initial interior point, and sparse workspaces
  \For{$outer=1$ to \texttt{max\_outer}}
    \For{$inner=1$ to \texttt{max\_inner}}
      \State Unpack \(x \rightarrow (v,\theta,v^{\qual{dc}},P_g,Q_g,P^{\qual{conv,ac}},Q^{\qual{conv,ac}},P^{\qual{conv,dc}})\)
      \State Evaluate equality constraints and Jacobian
      \State Evaluate objective gradient and diagonal Hessian terms
      \State Add bound barriers and nonlinear limit penalties
      \State Solve regularized KKT step
      \State Globalize by line search/restoration
      \If{acceptance fails} \State mark failure; \textbf{break} \EndIf
      \State Update interior point and penalties
      \If{primal/dual residuals below tolerance} \State converged; \textbf{break} \EndIf
    \EndFor
    \If{converged} \State \textbf{break} \EndIf
    \State Update barrier/penalty schedule
  \EndFor
  \If{failed and fallback allowed} \State call fast dispatch fallback \EndIf
  \State \Return OPF result
\EndProcedure
\end{algorithmic}

\paragraph{Algorithm A11: Parity NLP Build + Primal-Dual IPM}
\begin{algorithmic}[1]
\Procedure{\texttt{solve\_with\_parity\_ipm}}{$sys, opt$}
  \State Build full-space parity problem (AC, DC, converters, DER, shedding)
  \State Build bounds/scaling and initial point
  \State Initialize slacks/multipliers (\(z>0,\mu>0\))
  \For{$iter=1$ to \texttt{max\_iter}}
    \State Form condensed KKT (Mehrotra predictor-corrector)
    \State Solve affine predictor, then centering-corrector steps
    \State Optionally run Gondzio extra corrections
    \State Apply fraction-to-boundary step lengths
    \State Update \(x,\lambda,z,\mu\); re-evaluate residual metrics
    \State Track best iterate
    \If{primal, dual, and complementarity tolerances met}
      \State converged; \textbf{break}
    \EndIf
  \EndFor
  \State Restore best iterate when configured and needed
  \State Map \(x\) to \texttt{ACOPFResult} fields and return
\EndProcedure
\end{algorithmic}

\paragraph{Algorithm A12: Short-Circuit (All-Bus IEC-Style)}
\begin{algorithmic}[1]
\Procedure{\texttt{compute\_short\_circuit}}{$sys, sc\_opt$}
  \State Project to canonical AC model
  \State Build/factorize fault admittance matrix once
  \For{each faulted bus \(k\)}
    \State Solve \(Y^{\qual{fault}}z=e_k\), recover \(Z_{kk}\)
    \State Apply fault type and voltage-factor rule
    \State Compute \(I_{k}^{\prime\prime}\), short-circuit power, and \(R/X\)
    \State Append per-bus detailed result row
  \EndFor
  \State \Return short-circuit result set
\EndProcedure
\end{algorithmic}

\paragraph{Algorithm A13: Optimal Network Reconfiguration (ONR)}
\begin{algorithmic}[1]
\Procedure{\texttt{solve\_optimal\_reconfiguration}}{$sys, opt$}
  \State Project rich models to canonical AC network
  \State Build bus map, root selection, and net \(P/Q\) injections
  \State Build LinDistFlow MILP with \(\alpha,P,Q,v,f,t\) variables
  \State Solve MILP using native branch-and-cut + MIR cuts
  \State Round/repair \(\alpha\) to a radial spanning tree if needed
  \State Export open/closed stable branch indices and objective
  \State Run Newton PF on selected topology for verification
  \State \Return ONR result
\EndProcedure
\end{algorithmic}

\paragraph{Algorithm A14: Distribution PF (BFS)}
\begin{algorithmic}[1]
\Procedure{\texttt{solve\_distribution\_pf}}{$system, opt$}
  \If{$system$ is \texttt{HybridPowerSystem}}
    \State Project to canonical AC system
  \Else
    \State Use provided AC system directly
  \EndIf
  \State Build bus map and BFS tree from in-service branches
  \State Assemble net demand \(S^{\qual{net}}\)
  \State Initialize complex voltage vector
  \For{$iter=1$ to \texttt{max\_iter}}
    \State Backward sweep (currents: leaves \(\rightarrow\) root)
    \State Forward sweep (voltages: root \(\rightarrow\) leaves)
    \If{max voltage update \(< \texttt{tol}\)} \State converged; \textbf{break} \EndIf
  \EndFor
  \State Compute branch powers and total losses
  \State \Return DPF result
\EndProcedure
\end{algorithmic}

\paragraph{Algorithm A15: Carbon Analysis (Tracing + Matrix)}
\begin{algorithmic}[1]
\Procedure{\texttt{compute\_carbon\_analysis}}{$sys, pf, opt$}
  \State Require converged PF state
  \State Project to canonical hybrid model and build AC/DC id maps
  \State Build unified load sinks and explicit storage carbon-source states
  \State Build directed AC/DC/converter/router flow graph
  \State Build carbon sources and reconstruct slack contribution
  \State Run proportional tracing and export per-device allocations
  \State Solve hybrid matrix carbon-intensity system for bus intensities
  \State Export summaries, balance checks, and solver-specific notes
  \State \Return carbon-analysis result
\EndProcedure
\end{algorithmic}

\paragraph{Algorithm A16: DC Optimal Power Flow}
\begin{algorithmic}[1]
\Procedure{\texttt{solve\_dc\_opf}}{$sys, opt$}
  \State Build bus/branch/generator index sets and susceptance matrix $B$
  \State Assemble cost vector (linear or quadratic depending on backend)
  \State Add per-bus load-shedding variables with auto-tuned VOLL penalty
  \State Build power-balance equality constraints (per bus)
  \State Add flow-definition equalities: $P_{f,ij} - b_{ij}(\theta_i - \theta_j) = 0$
  \State Set variable bounds ($\theta$, $P_g$, $P_f$, $\Delta P_d$)
  \State Select backend: NativeQP $\to$ Gurobi $\to$ HiGHS $\to$ Dual Simplex
  \State Solve LP or QP
  \State Extract LMPs from power-balance dual variables
  \State \Return DC OPF result with dispatch, flows, and LMPs
\EndProcedure
\end{algorithmic}

\paragraph{Algorithm A17: Reactive Power Optimization (OA)}
\begin{algorithmic}[1]
\Procedure{\texttt{solve\_reactive\_power\_opt}}{$sys, opt$}
  \State Enumerate discrete variables (transformer taps, switchable shunt steps)
  \State Solve baseline AC OPF to obtain incumbent objective
  \State \textbf{Phase 0 (Sensitivity):} For each discrete variable, perturb $\pm 1$, solve NLP, rank by $|\Delta\text{obj}|$
  \State \textbf{Phase 1 (Ternary search):} For each variable in sensitivity order, run integer ternary search; add OA cuts $f(y) \ge f_k + g_k^T(y - y_k)$ at each evaluation
  \State \textbf{Phase 2 (Coordinate descent, up to 3 cycles):} Sweep each variable; skip candidates where OA lower bound $\ge$ incumbent; add cuts; terminate if no improvement
  \State \textbf{Phase 3 (Pairwise polishing):} Check $\pm 1$ pairs of discrete variables; OA prunes before NLP evaluation
  \State \Return best tap/shunt settings, OA cuts, NLP solve count
\EndProcedure
\end{algorithmic}

\paragraph{Algorithm A18: LCQP Interior-Point Solver}
\begin{algorithmic}[1]
\Procedure{\texttt{solve\_lcqp}}{$Q, c, A_{eq}, b_{eq}, \ell, u$}
  \State Initialize $x$ at bound midpoints; $s_\ell, s_u \gets$ slacks; $z_\ell, z_u \gets 1$; $y \gets 0$
  \State Cache symbolic factorization of KKT sparsity pattern
  \For{$iter = 1$ to \texttt{max\_iter}}
    \State Compute complementarity gap $\mu = (s_\ell^T z_\ell + s_u^T z_u) / (2n_{\text{bounded}})$
    \If{primal, dual feasibility, and gap $< \epsilon$} \State converged; \textbf{break} \EndIf
    \State Form KKT residuals (stationarity, primal feasibility, complementarity)
    \State Solve predictor step via sparse LU with $\epsilon$-regularization
    \State Compute centering parameter $\sigma = \min(0.3, 100\mu)$
    \State Solve corrector step with second-order compensation
    \State Compute step lengths $\alpha_p, \alpha_d$ via fraction-to-boundary ($\tau = 0.995$)
    \State Update $(x, s_\ell, s_u, y, z_\ell, z_u)$
  \EndFor
  \State \Return $x$, objective, convergence status
\EndProcedure
\end{algorithmic}

\paragraph{Algorithm A19: Dual Simplex LP Solver}
\begin{algorithmic}[1]
\Procedure{\texttt{solve\_dual\_simplex}}{$c, A, b, \ell, u$}
  \State Convert to standard form: add slacks, surplus, artificial variables
  \State \textbf{Crash basis:} For each artificial, find original column with largest $|A_{ij}|$ and swap
  \State \textbf{Phase I:} Minimize sum of artificials
  \For{$iter = 1$ to \texttt{max\_iter}}
    \State Compute reduced costs; select entering variable (most-negative)
    \State Compute step direction and leaving variable (ratio test)
    \If{all artificials non-basic} \State Phase I complete; \textbf{break} \EndIf
    \State Pivot basis; refactorize every 15--50 iterations
  \EndFor
  \If{artificials remain in basis} \State Extract Farkas certificate; \Return infeasible \EndIf
  \State \textbf{Phase II:} Minimize original $c^T x$
  \For{$iter = 1$ to \texttt{max\_iter}}
    \State Compute reduced costs
    \If{all non-negative (dual feasible)} \State optimal; \textbf{break} \EndIf
    \State Select entering/leaving variables and pivot
  \EndFor
  \State \Return $x$, objective, basis (cached for warm-start in branch-and-cut)
\EndProcedure
\end{algorithmic}

\paragraph{Algorithm A20: Annual Production Simulation}
\begin{algorithmic}[1]
\Procedure{\texttt{solve\_annual\_production\_simulation}}{$sys, ts\_data, opts$}
  \State \textbf{L0 (Annual planning):} Build coarse monthly/weekly blocks
  \Statex \quad Generate energy budgets $\bar E_{g,k}$, SOC trajectory, maintenance flags
  \For{each block $k$}
    \State \textbf{L2 (Weekly UC):} Extract sub-horizon and solve rolling UC
    \Statex \quad Binding window: $T^{(2)}$ steps; look-ahead: $T^{\qual{la}}$ steps
    \Statex \quad Call \texttt{solve\_unit\_commitment(sub\_sys, sub\_ts, ...)}
    \If{not \texttt{skip\_replay}}
      \For{each day in block}
        \State \textbf{L3 (Daily replay):} Slice UC schedule for the day
        \State Call \texttt{solve\_time\_series\_pf(sub\_sys, ...)} for OPF/PF validation
        \State Store per-step voltages, losses, flows, costs
      \EndFor
    \Else
      \State Fill step results from UC + load profiles (no OPF/PF)
    \EndIf
  \EndFor
  \State Aggregate generator/renewable/storage annual statistics
  \Statex \quad Capacity factors, curtailment rates, cycle counts, startups
  \State \Return annual result with step/block/component summaries
\EndProcedure
\end{algorithmic}

\paragraph{Algorithm A21: Lifecycle Simulation}
\begin{algorithmic}[1]
\Procedure{\texttt{run\_lifecycle\_simulation}}{$sys, opts$}
  \For{$year = 1$ to \texttt{num\_years}}
    \State Apply load growth: $P_d^{(y)} = P_d^{(0)} (1 + r_g)^{y-1}$
    \State Apply PV derating: $P_{\max}^{(y)} = P_{\max}^{(0)} (1 - d_{PV})^{y-1}$
    \State Compute battery SOH: $\min(\text{SOH}_{\text{cal}}, \text{SOH}_{\text{cyc}})$
    \If{$\text{SOH} \le \text{SOH}_{\text{EOL}}$}
      \State Trigger replacement; reset SOH and accumulate replacement cost
    \EndIf
    \State \textbf{Tier 1:} Run annual chronological dispatch
    \State \textbf{Tier 2:} Stratified sampling into 7 operating regimes
    \Statex \quad (peak load, low renewable, high curtailment, heavy charge/discharge, normal, low load)
    \State Compute Neyman-allocated sampling error bounds
    \State Compute dispatch, sampling, and storage carbon error bounds
    \State Accumulate annual OpEx and carbon metrics
  \EndFor
  \State Compute NPV: $\sum_{y} (\text{OpEx}_y + \text{ReplaceEx}_y) / (1+r)^{y-1}$
  \State \Return lifecycle result with per-year metrics and error bounds
\EndProcedure
\end{algorithmic}
